\label{der:6-4}
规范形中正坐标空间给出$p$维正定子空间。任意维数大于$p$的子空间，与维数$n-p$的非正坐标空间维数和大于$n$，有非零交向量，在该向量上二次型不可能严格正。因此最大正定子空间维数恰为$p$；对$-A$同样得到最大负定维数为$q$。

若$A=Q\Lambda Q\T\succ0$，定义$S=Q\diag(\sqrt{\lambda_i})Q\T$。它实对称，谱全正，故正定；相乘时$Q\T Q=I$，所以$S^2=Q\Lambda Q\T=A$。若另有正定对称$T$满足$T^2=A$，则$TA=T^3=AT$，故$T$保留$A$的每个特征空间。对$A$的$\lambda$特征空间，$T$的限制仍正定对称且平方为$\lambda I$；正交对角化该限制，所有特征值$\mu$都满足$\mu^2=\lambda,\mu>0$，只能为$\sqrt\lambda$，所以限制恰为$\sqrt\lambda I$。各特征空间直和为全空间，故$T=S$，证明唯一。

最后，谱分解写为
\[A=\sum_{\lambda_i\ne0}\sgn(\lambda_i)\bigl(\sqrt{|\lambda_i|}q_i\bigr)\bigl(\sqrt{|\lambda_i|}q_i\bigr)\T.\]
若全用正号，任意$x$代入后都是平方和，必半正定；反向半正定时全部非零谱值正，上式全为正号。由此外积表示的准确条件得到证明。
